\documentclass[a4paper,14pt, unknownkeysallowed]{extreport}
\usepackage{placeins}
\usepackage{cmap} % Улучшенный поиск русских слов в полученном pdf-файле
\usepackage[T2A]{fontenc} % Поддержка русских букв
\usepackage[utf8]{inputenc} % Кодировка utf8
\usepackage[english,russian]{babel} % Языки: русский, английский
%\usepackage{pscyr} % Нормальные шрифты
\usepackage{enumitem}
\usepackage{multirow}
\usepackage[14pt]{extsizes}

\usepackage{caption}
\captionsetup{labelsep=endash}
\captionsetup[figure]{name={Рисунок}}

\usepackage{amsmath}

\usepackage{geometry}
\geometry{left=30mm}
\geometry{right=15mm}
\geometry{top=20mm}
\geometry{bottom=20mm}

\setlength{\parindent}{1.25cm} % для шрифта 14pt

\usepackage{titlesec}
\titleformat{\section}
	{\normalsize\bfseries}
	{\thesection}
	{1em}{}
\titlespacing*{\chapter}{0pt}{-30pt}{8pt}
\titlespacing*{\section}{\parindent}{*4}{*4}
\titlespacing*{\subsection}{\parindent}{*4}{*4}

\usepackage{setspace}
\onehalfspacing % Полуторный интервал

\frenchspacing
\usepackage{indentfirst} % Красная строка

\usepackage{titlesec}
\titleformat{\chapter}{\LARGE\bfseries}{\thechapter}{20pt}{\LARGE\bfseries}
\titleformat{\section}{\Large\bfseries}{\thesection}{20pt}{\Large\bfseries}

\usepackage{listings}
\usepackage{xcolor}

% Для листинга кода:
\lstset{ %
	language=lisp,   					% выбор языка для подсветки	
	basicstyle=\small\sffamily,			% размер и начертание шрифта для подсветки кода
	numbers=left,						% где поставить нумерацию строк (слева\справа)
	%numberstyle=,					% размер шрифта для номеров строк
	stepnumber=1,						% размер шага между двумя номерами строк
	numbersep=5pt,						% как далеко отстоят номера строк от подсвечиваемого кода
	frame=single,						% рисовать рамку вокруг кода
	tabsize=4,							% размер табуляции по умолчанию равен 4 пробелам
	captionpos=t,						% позиция заголовка вверху [t] или внизу [b]
	breaklines=true,					
	breakatwhitespace=true,				% переносить строки только если есть пробел
	escapeinside={\#*}{*)},				% если нужно добавить комментарии в коде
	backgroundcolor=\color{white},
}
\usepackage{pgfplots}
\usetikzlibrary{datavisualization}
\usetikzlibrary{datavisualization.formats.functions}

\usepackage{graphicx}
\newcommand{\img}[3] {
	\begin{figure}[h!]
		\center{\includegraphics[height=#1]{inc/img/#2}}
		\caption{#3}
		\label{img:#2}
	\end{figure}
}
\newcommand{\boximg}[3] {
	\begin{figure}[h]
		\center{\fbox{\includegraphics[height=#1]{inc/img/#2}}}
		\caption{#3}
		\label{img:#2}
	\end{figure}
}

% для матриц
\makeatletter
\renewcommand*\env@matrix[1][\arraystretch]{%
	\edef\arraystretch{#1}%
	\hskip -\arraycolsep
	\let\@ifnextchar\new@ifnextchar
	\array{*\c@MaxMatrixCols c}}
\makeatother

\usepackage{amsfonts}
\usepackage{graphicx}


\usepackage[justification=centering]{caption} % Настройка подписей float объектов

\usepackage[unicode,pdftex]{hyperref} % Ссылки в pdf
\hypersetup{hidelinks}

\usepackage{csvsimple}

\addto\captionsrussian{\def\refname{Список используемых источников}}

\newcommand{\code}[1]{\texttt{#1}}

\makeatletter
\renewcommand{\@biblabel}[1]{#1. \hfill}
\makeatother


\begin{document}


\begin{titlepage}
	\newgeometry{pdftex, left=2cm, right=2cm, top=2.5cm, bottom=2.5cm}
	\fontsize{12pt}{12pt}\selectfont
	\noindent \begin{minipage}{0.15\textwidth}
		\includegraphics[width=\linewidth]{inc/img/b_logo.jpg}
	\end{minipage}
	\noindent\begin{minipage}{0.9\textwidth}\centering
		\textbf{Министерство науки и высшего образования Российской Федерации}\\
		\textbf{Федеральное государственное бюджетное образовательное учреждение высшего образования}\\
		\textbf{«Московский государственный технический университет имени Н.Э.~Баумана}\\
		\textbf{(национальный исследовательский университет)»}\\
		\textbf{(МГТУ им. Н.Э.~Баумана)}
	\end{minipage}
	
	\noindent\rule{18cm}{3pt}
	\newline\newline
	\noindent ФАКУЛЬТЕТ $\underline{\text{«Информатика и системы управления»~~~~~~~~~~~~~~~~~~~~~~~~~~~~~~~~~~~~~~~~~~~~~~~~~~~~~~~}}$ \newline\newline
	\noindent КАФЕДРА $\underline{\text{«Программное обеспечение ЭВМ и информационные технологии»~~~~~~~~~~~~~~~~~~~~~~~}}$\newline\newline\newline\newline\newline\newline\newline
	
	
	\begin{center}
		\Large\textbf{Отчет по лабораторной работе №2}
		\Large\textbf {по дисциплине <<Функциональное и логическое программирование>>}
	\end{center}
	
	\noindent\textbf{Тема} $\underline{\text{~~Определений функций пользователя~~~~~~~~~~~~~~~~~~~~~~~~~~~~~~~~~~~~~~~~~~~~~~~~~~~~~~~~~~~~~~~~~}}$\newline\newline\newline
	\noindent\textbf{Студент} $\underline{\text{~~Котляров Н.А.~~~~~~~~~~~~~~~~~~~~~~~~~~~~~~~~~~~~~~~~~~~~~~~~~~~~~~~~~~~~~~~~~~~~~~~~~~~~~~~~~~~~~~~~~}}$\newline\newline
	\noindent\textbf{Группа} $\underline{\text{~~ИУ7-61Б~~~~~~~~~~~~~~~~~~~~~~~~~~~~~~~~~~~~~~~~~~~~~~~~~~~~~~~~~~~~~~~~~~~~~~~~~~~~~~~~~~~~~~~~~~~~~~~~~}}$\newline\newline
	\noindent\textbf{Оценка (баллы)} $\underline{\text{~~~~~~~~~~~~~~~~~~~~~~~~~~~~~~~~~~~~~~~~~~~~~~~~~~~~~~~~~~~~~~~~~~~~~~~~~~~~~~~~~~~~~~~~~~~~~~~~~~~}}$\newline\newline
	\noindent\textbf{Преподаватель} $\underline{\text{~~Толпинская Н.~Б.~~Строганов~Ю.~В.~~~~~~~~~~~~~~~~~~~~~~~~~~~~~~~~~~~~~~~~~~~~~~~~~~}}$\newline
	
	\begin{center}
		\vfill
		Москва~---~\the\year
		~г.
	\end{center}
	\restoregeometry
\end{titlepage}



\section*{Задание 1}
Создать базу знаний «Предки», позволяющую наиболее эффективным способом
(за меньшее количество шагов, что обеспечивается меньшим количеством предложений БЗ - правил), и используя разные варианты (примеры) простого вопроса определить:
\begin{enumerate}
	\item по имени субъекта определить всех его бабушек,
	\item по имени субъекта определить всех его дедушек,
	\item по имени субъекта определить всех его бабушек и дедушек,
	\item по имени субъекта определить его бабушку по материнской линии,
	\item по имени субъекта определить его бабушку и дедушку по материнской линии.
\end{enumerate}

\begin{lstlisting}
parent(olga, oleg, f).
parent(ivan, oleg, m).
parent(oleg, yaroslav, m).
parent(inga, yaroslav, f).
parent(oleg, kate, m).
parent(inga, kate, f).
parent(nikita, georg, m).
parent(sveta, georg, f).
parent(yaroslav, stepan, m).
parent(irina, stepan, f).
parent(georg, sergey, m).
parent(kate, sergey, f).
parent(georg, alex, m).
parent(kate, alex, f).
parent(georg, nika, m).
parent(kate, nika, f).
grandparent(X, Y)   :- parent(X, Z, _), parent(Z, Y, _).
grandmother(X, Y)   :- parent(X, Z, f), parent(Z, Y, _).
grandfather(X, Y)   :- parent(X, Z, m), parent(Z, Y, _).
maternal_grandmother(X, Y) :- parent(X, Z, f), parent(Z, Y, f).
maternal_grandparent(X, Y) :- parent(X, Z, _), parent(Z, Y, f).	
\end{lstlisting}

\section*{Задание 2}

Дополнить базу знаний правилами, позволяющими найти
\begin{enumerate}
	\item Максимум из двух чисел
	\begin{enumerate}
		\item без использования отсечения,
		\item с использованием отсечения;
	\end{enumerate}
	\item Максимум из трех чисел
	\begin{enumerate}
		\item без использования отсечения,
		\item с использованием отсечения;
	\end{enumerate}
\end{enumerate}

\begin{lstlisting}
max2(N1, N2, N2)    :- N2 >= N1.
max2(N1, N2, N1)    :- N1 >= N2.

max3(N1, N2, N3, N1)    :- N1 >= N2, N1 >= N3.
max3(N1, N2, N3, N2)    :- N2 >= N1, N2 >= N3.
max3(N1, N2, N3, N3)    :- N3 >= N1, N3 >= N2.

max2clip(N1, N2, N2)    :- N2 >= N1, !.
max2clip(N1, _, N1).

max3clip(N1, N2, N3, N1)   :- N1 >= N2, N1 >= N3, !.
max3clip(_, N2, N3, N2)    :- N2 >= N3, !.
max3clip(_, _, N3, N3).
\end{lstlisting}

\begin{table}[ht!]
	\centering
	\label{decisions}
	\begin{tabular}{|p{1.1cm}|p{5cm}|p{6cm}|p{5cm}|}
		\hline
		\textbf{Шаг} & \textbf{Состояние резольвенты} & \textbf{Для каких термов запускается алгоритм унификации и результат подстановки} & \textbf{Дальнейшие действия: прямой ход или откат}\\
		\hline
		1 & max3(3, 2, 1, Asw) & max3(3, 2, 1, Asw) = max2(N1, N2, N2) Подстановка: нет & Откат \\
		\hline
		... & & & \\
		\hline
		3 & 3 >= 2, 3 >= 1 & max3(3, 2, 1, Asw) = max3(N1, N2, N3, N1) Подстановка:\{N1=3; N2=2; N3=1; Asw=3\} & Прямой ход \\
		\hline
		4 & 3 >= 1 & 3 >= 2 Верно. Подстановка:\{N1=3; N2=2; N3=1; Asw=3\} & Прямой ход \\
		\hline
		5 & Пусто & 3 >=1 Верно. Подстановка:\{N1=3; N2=2; N3=1; Asw=3\} & Конкретизация Asw. Откат к пункту 3. Дальнейшее продолжение перебора. \\
		\hline
	\end{tabular}
\end{table}

\begin{table}[ht!]
	\centering
	\label{decisions}
	\begin{tabular}{|p{1.1cm}|p{5cm}|p{6cm}|p{5cm}|}
		\hline
		\textbf{Шаг} & \textbf{Состояние резольвенты} & \textbf{Для каких термов запускается алгоритм унификации и результат подстановки} & \textbf{Дальнейшие действия: прямой ход или откат}\\
		\hline
		1 & max3clip(2, 3, 1, Asw) & max3clip(2, 3, 1, Asw) = max2(N1, N2, N2) Подстановка: нет & Откат \\
		\hline
		... & & & \\
		\hline
		8 & 2 >= 3, 2 >= 1 & max3clip(2, 3, 1, Asw) = max3clip(N1, N2, N3, N1) Подстановка:\{N1=2; N2=3; N3=1; Asw=2\} & Прямой ход \\
		\hline
		9 & 2 >= 1 & 2 >= 3 Неверно. & Откат к пункту 8 \\
		\hline
		10 & 3 >= 1 & max3clip(2, 3, 1, Asw) = max3clip(\_, N2, N3, N2) Подстановка: \{N2=3; N3=1; Asw=3\} & Прямой ход \\
		\hline
		11 & Пусто & 3 >= 1 Верно. Подстановка: \{N2=3; N3=1; Asw=3\} & Конкретизация Asw. Оператор отсечения. Откат к пункту 8. Завершение работы.\\
		\hline
	\end{tabular}
\end{table}

\end{document}